\chapter{Experiment design: the manoeuvres}
\label{ch:maneuvers}

The quality of an identified model is limited by the data it is fitted to. No estimator can
recover a parameter that the experiment did not excite. This chapter covers the general
principles and then each manoeuvre in \file{maneuvers.md}.

\section{Principles}

\paragraph{Excitation.} Least-squares estimation (\cref{ch:estimation}) needs the regressors, here
$\dot u$, $u$ and $|u|u$, to vary \emph{independently} over the data set. If the vehicle only
ever cruises at one speed, $u$ and $|u|u$ are proportional and $X_u$ cannot be separated
from $\Xuu$. If it never accelerates, $\mtot$ is invisible. This property is called
\emph{persistent excitation} \cite{ljung1999}. Different manoeuvres excite different
regressors:
\begin{itemize}[nosep]
  \item steps from rest: large $\dot u$ at small $u$, which isolates $\mtot$;
  \item steady states at several levels: $\dot u = 0$ at different $u$, which separates $X_u$ from $\Xuu$;
  \item coast-down: $T = 0$, so no thrust error enters;
  \item sinusoids: controlled ratio of $\dot u$ to $u$ through the frequency.
\end{itemize}

\paragraph{Open loop.} The command scripts drive \code{/pixhawk/manual_control} in
MANUAL mode, never \code{/pixhawk/cmd_vel} in GUIDED mode. In GUIDED mode ArduSub's velocity
controller would close a loop around the vehicle. The thrust would then react to the
measured velocity, including its disturbances and noise, and input and noise become
correlated. Least squares assumes they are not, and the estimates become biased: closed-loop
data tends to identify the inverse of the controller rather than the plant
\cite[ch.~13]{ljung1999}. The autonomy bags from Zermatt are closed loop for the same reason,
which makes them acceptable for validation but not ideal for fitting.

\paragraph{Repeatability and holdout.} Every manoeuvre is repeated, and runs are split into a
fitting set and a held-out validation set, so that the validation answers ``does the model
predict data it has not seen?''.

\section{Command interface}

\begin{table}[htbp]
\centering
\caption{Axes of the \code{/pixhawk/manual_control} message (\code{Int16MultiArray}).}
\label{tab:axes}
\begin{tabular}{@{}c l l c@{}}
\toprule
Index & Axis & Range & Neutral \\
\midrule
0 & $x$ surge & $-1000 \dots 1000$ & 0 \\
1 & $y$ sway  & $-1000 \dots 1000$ & 0 \\
2 & $z$ heave & $0 \dots 1000$ & 500 \\
3 & $r$ yaw   & $-1000 \dots 1000$ & 0 \\
4 & $s$ roll  & $-1000 \dots 1000$ & 0 \\
5 & $t$ pitch & $-1000 \dots 1000$ & 0 \\
\bottomrule
\end{tabular}
\end{table}
\FloatBarrier

\code{_step_common.run_constant_surge()} publishes \code{[surge, 0, 500, 0, 0, 0]} at
\qty{20}{\hertz} for the hold time, clamps surge to $\pm500$, and then sends ten neutral
messages so that a single lost packet cannot leave the thruster running. The vehicle has to
be armed and in MANUAL mode beforehand; the script deliberately does neither.

\paragraph{Event markers.} Each run publishes JSON markers on \code{/sysid/event}:
\code{thrust_on} with the first command, \code{thrust_cut} with the first neutral message,
and \code{done} (or \code{interrupted}). The fit scripts can therefore cut a bag into
segments without guessing times from the PWM, and every run provides two sharp edges for
the delay measurement of \cref{ch:timing}. The markers only help if \code{/sysid/event} is
in the recorder's topic list.

\section{Static tests (pool)}

\subsection{Static tilt}

A known moment $M_\mathrm{app}$ (a weight at a known lever) is applied and the steady tilt is
measured. From \eqref{eq:pitch-stiffness},
\begin{equation}
  M_\theta = \frac{M_\mathrm{app}}{\theta - \theta_0}.
\end{equation}
Several weights give several points and a straight-line fit. Curvature would reveal
large-angle nonlinearity ($\sin\theta$ rather than $\theta$).

\subsection{Free decay}

The vehicle is displaced in pitch (or roll) and released. The rotational equation
\begin{equation}
  (I_{yy} - M_{\dot q})\,\ddot\theta + M_q\,\dot\theta + M_{q|q|}|\dot\theta|\dot\theta + M_\theta\,\theta = 0
  \label{eq:free-decay}
\end{equation}
is a damped oscillator. In its linear form the standard quantities are
\begin{equation}
  \omega_n = \sqrt{\frac{M_\theta}{I_{yy} - M_{\dot q}}}, \qquad
  \zeta = \frac{M_q}{2\sqrt{M_\theta\,(I_{yy} - M_{\dot q})}}, \qquad
  \omega_d = \omega_n\sqrt{1-\zeta^2}.
\end{equation}
With $M_\theta$ known from the tilt test, the measured frequency gives the total inertia
$I_{yy} - M_{\dot q} = M_\theta/\omega_n^2$. The decay of successive peaks gives the damping
through the logarithmic decrement,
\begin{equation}
  \delta = \ln\frac{\theta_k}{\theta_{k+1}} = \frac{2\pi\zeta}{\sqrt{1-\zeta^2}} .
\end{equation}
With quadratic damping the decrement depends on amplitude: large swings decay
proportionally faster. Plotting $\delta$ against amplitude separates the linear from the
quadratic part. The equivalent linear damping of a quadratic term at amplitude $A$ and
frequency $\omega$ is $\tfrac{8}{3\pi}M_{q|q|}\,\omega A$ \cite[sec.~6.4]{fossen2011}.

\begin{figure}[htbp]
\centering
\begin{tikzpicture}
\begin{axis}[wide,xlabel={$t$ / \unit{\second}},ylabel={$\theta$ / \unit{\radian}},xmin=0,xmax=8]
\addplot[black] table[x=t,y=dummy]{data/decay.dat};
\addplot[blue,dashed] table[x=t,y=lowdamp]{data/decay.dat};
\legend{dummy coefficients ($\zeta \approx 0.58$),linear damping $\div 10$}
\end{axis}
\end{tikzpicture}
\caption{Simulated pitch free decay from \qty{0.35}{\radian} with the dummy BlueROV2
coefficients ($I_{yy} - M_{\dot q} = \qty{0.365}{\kilogram\meter\squared}$,
$M_\theta = \qty{1.316}{\newton\meter\per\radian}$, $M_q = \qty{0.8}{\newton\meter\second\per\radian}$,
$M_{q|q|} = \qty{0.47}{\newton\meter\second\squared\per\radian\squared}$) and with the linear
damping reduced tenfold.}
\label{fig:decay}
\end{figure}
\FloatBarrier

\Cref{fig:decay} illustrates a practical limitation. With the dummy coefficients pitch is so
heavily damped ($\zeta\approx0.58$) that barely one overshoot is visible. Frequency and
decrement are then poorly determined from peaks. If POLARIS turns out similar, the decay
should be fitted as a whole trajectory by integrating \eqref{eq:free-decay}
(output error, \cref{sec:output-error}) rather than read off peak by peak.

\section{Dynamic tests (lake)}

\subsection{Step from rest}

At $t = 0^+$ the vehicle is at rest, so all damping vanishes and \eqref{eq:surge} gives
$\dot u(0^+) = T/\mtot$. The initial acceleration therefore measures the total inertia
directly. \Cref{ch:timing} explains why this is harder than it looks.

\subsection{Multi-level steps}

The four scripts \code{step_surge_200} to \code{step_surge_500} hold a constant command for
\qty{10}{\second}. Averaging the speed over the steady part of each run gives pairs
$(T_i, u_{\mathrm{ss},i})$ that satisfy
\begin{equation}
  T_i = X_u\,u_{\mathrm{ss},i} + \Xuu\,u_{\mathrm{ss},i}^2 .
  \label{eq:ss-regression}
\end{equation}
This is \emph{linear in the unknowns} $X_u$ and $\Xuu$ even though the model is nonlinear in
$u$, so ordinary least squares applies directly. Two levels determine the two unknowns, and
more levels give a check that the quadratic law actually fits. The levels should span the
operating speed range, because the low-speed points carry most of the information about $X_u$.

\subsection{Coast-down}

\code{coastdown_surge_400.py} holds 400 for \qty{20}{\second}, which is long enough to reach steady speed and to
average DVL noise, and then cuts thrust at the \code{thrust_cut} marker. The decay after the cut follows
\eqref{eq:decay-both} with no thrust term, so no thrust-curve error enters. The
linearisations of \cref{sec:closedform} give quick checks: $\ln u$ against $t$ is straight
for linear damping and $1/u$ against $t$ for quadratic damping (\cref{fig:coast-lin}). In
practice \eqref{eq:decay-both} is fitted directly, and the result is the pair of ratios
$X_u/\mtot$ and $\Xuu/\mtot$ (\cref{sec:identifiability}).

\begin{figure}[htbp]
\centering
\begin{subfigure}{0.48\textwidth}
\centering
\begin{tikzpicture}
\begin{axis}[notes,xlabel={$t$ / \unit{\second}},ylabel={$1/u$ / \unit{\second\per\meter}},xmin=0,xmax=6]
\addplot[red] table[x=t,y=invquad]{data/coast.dat};
\end{axis}
\end{tikzpicture}
\caption{Quadratic only: $1/u$ against $t$}
\end{subfigure}\hfill
\begin{subfigure}{0.48\textwidth}
\centering
\begin{tikzpicture}
\begin{axis}[notes,xlabel={$t$ / \unit{\second}},ylabel={$\ln(u / \unit{\meter\per\second})$},xmin=0,xmax=6]
\addplot[blue] table[x=t,y=lnlin]{data/coast.dat};
\end{axis}
\end{tikzpicture}
\caption{Linear only: $\ln u$ against $t$}
\end{subfigure}
\caption{Linearising transformations of the simulated coast-down curves of
\cref{fig:step}: reciprocal speed for purely quadratic damping and logarithm of speed for purely
linear damping, both against time.}
\label{fig:coast-lin}
\end{figure}
\FloatBarrier

\subsection{Forced oscillation}

A sinusoidal command at frequency $\omega$ keeps the vehicle within a short travel distance.
For the linearised surge dynamics $\mtot\dot u + d\,u = F$ the frequency response is
\begin{equation}
  G(j\omega) = \frac{U}{F} = \frac{1}{d + j\omega\mtot}
  \quad\Rightarrow\quad
  \mathrm{Re}\,G^{-1} = d, \qquad \frac{\mathrm{Im}\,G^{-1}}{\omega} = \mtot .
\end{equation}
Amplitude and phase at each frequency therefore give damping and inertia \emph{simultaneously}.
Low frequencies are dominated by damping and high frequencies by inertia, which is why the
sweep spans a range. With quadratic damping, $d$ is the amplitude-dependent equivalent
damping $X_u + \tfrac{8}{3\pi}\Xuu U$.

\subsection{Combined axes}

Surge and yaw commanded together excite the products $ur$ and $vr$ that single-axis
manoeuvres never produce (\cref{ch:symmetry}). These runs test the coupling, see \cref{ch:coupling}.

\section{Summary}

\begin{table}[htbp]
\centering
\caption{Manoeuvres of \file{maneuvers.md}, the quantities each one determines, and the
implementation state of the corresponding scripts.}
\label{tab:maneuvers}
\small
\begin{tabularx}{\textwidth}{@{}l X l@{}}
\toprule
Manoeuvre & Determines & Scripts \\
\midrule
Static tilt & $M_\theta$ and $\theta_0$ (roll part dropped 15 Oct) & fit placeholder \\
Free decay & $I_{yy} - M_{\dot q}$; pitch damping (roll part dropped 15 Oct) & fit placeholder \\
Step from rest & $\mtot$ (from $T/\dot u(0^+)$) & command scripts exist \\
Multi-level steps & $X_u$, $\Xuu$ (in thrust units) & \code{step_surge_200..500} \\
Coast-down & $X_u/\mtot$, $\Xuu/\mtot$ & \code{coastdown_surge_400} \\
Forced oscillation & $\mtot$ and damping together & not yet \\
Combined axes & coupling terms & not yet \\
\bottomrule
\end{tabularx}
\end{table}
\FloatBarrier

\begin{revised}[title=Revised 15 October 2026]
Roll is no longer modelled (\cref{sec:no-roll}). The roll parts of the static tilt and the roll
free decay (section 2.1.3 of \file{maneuvers.md}) are not needed, and \cref{tab:maneuvers} lists
only the pitch quantities. The theory of this chapter applies to roll unchanged if roll is
ever re-enabled.
\end{revised}

\begin{questions}
\begin{enumerate}
  \item Why can a single steady-state run not separate $X_u$ from $\Xuu$?
  \item The pitch free decay of a vehicle shows peaks of \num{0.30}, \num{0.12} and
        \qty{0.05}{\radian}. Estimate $\zeta$ from the first decrement. Does the damping look linear?
  \item Why does a coast-down not need the thrust curve, and what can it still not deliver on its own?
\end{enumerate}
\end{questions}
